\section{Absolute scores}
\label{app:absolute}
The main text reports differences, because the level of a Brier score is set by
how hard a window's questions happened to be. Table~\ref{tab:absolute} gives
the levels themselves, so readers need not reconstruct them. Uncertainty (the
score of a forecaster who answers the window's base rate on every row) is a
property of the outcomes, so it is the same for everyone in a window and appears
in the heading. Skill is the share of that forecaster's error removed.

The constant forecasters are included to make the scale legible, not because
they are contenders. Answering Yes on every row scores worse than the
base rate by a factor of two or more, and answering 0.5 on every row scores
worse than the base rate on all three windows. Beating them is not evidence of
judgment; the informative comparisons are with the market and the shrunk market.

The shrunk market is a linear, and so monotone, compression of the price. Applied
to the price it is pure calibration loss, yet it reproduces 24 to 52 per cent of
the models' shortfall while calibration explains much less of it. The two fit
together: a calibrated forecaster that knows less than the market also varies
less than the price, and in that forecaster the compression is lost resolution,
not miscalibration.

\begin{table*}[!tbp]
\centering
\caption{Absolute scores on the test split of the primary set. Brier and
reliability are lower-is-better; resolution and skill are higher-is-better. Constant
forecasters have no resolution by construction and their reliability is not comparable
with a varying forecaster's, so those cells are left blank. Reliability minus resolution plus
uncertainty differs from Brier by up to 0.002, the within-bin term that ten bins leave
out~\citep{stephenson2008}. The market rows average the market's score over every model's
rows stacked together, so they differ in the fourth decimal from the market score behind any
one model's Brier edge, which is computed row by row on that model's own rows. The Brier edge (market minus model, on each model's own rows) and the
market-minus-model resolution carry 95\% event-resampling ranges.}
\label{tab:absolute}
\begingroup\footnotesize\setlength{\tabcolsep}{5pt}%
\begin{tabular}{@{}lrrlrrlr@{}}
\toprule
Forecaster & Rows & Brier & Brier edge & REL & RES & Market minus model RES & Skill \\
\midrule
\multicolumn{8}{@{}l}{\textbf{\winA} (uncertainty 0.2317, Yes base rate 0.365)}\\[1pt]
\mGptFiveTwo & 1{,}248 & 0.2007 & $-$0.0153 {\scriptsize[$-$0.0265, $-$0.0051]} & 0.0051 & 0.0352 & +0.013 {\scriptsize[+0.003, +0.024]} & +0.134 \\
\mGptOss & 1{,}248 & 0.2159 & $-$0.0305 {\scriptsize[$-$0.0449, $-$0.0162]} & 0.0074 & 0.0214 & +0.027 {\scriptsize[+0.015, +0.040]} & +0.068 \\
\mKimiKTwoFive & 1{,}237 & 0.1990 & $-$0.0129 {\scriptsize[$-$0.0264, $-$0.0013]} & 0.0033 & 0.0352 & +0.013 {\scriptsize[+0.002, +0.024]} & +0.142 \\
\mDeepSeekVThreeTwo & 1{,}248 & 0.1997 & $-$0.0143 {\scriptsize[$-$0.0262, $-$0.0030]} & 0.0023 & 0.0338 & +0.015 {\scriptsize[+0.003, +0.028]} & +0.138 \\
\addlinespace
\market & --- & 0.1855 & --- & 0.0015 & 0.0486 & --- & +0.199 \\
Market, shrunk & --- & 0.1951 & --- & --- & --- & --- & +0.158 \\
Base rate every row & --- & 0.2317 & --- & --- & 0.0000 & --- & +0.000 \\
0.5 every row & --- & 0.2500 & --- & --- & 0.0000 & --- & $-$0.079 \\
Yes every row & --- & 0.6352 & --- & --- & 0.0000 & --- & $-$1.741 \\
No every row & --- & 0.3648 & --- & --- & 0.0000 & --- & $-$0.574 \\
\addlinespace
\multicolumn{8}{@{}l}{\textbf{\winB} (uncertainty 0.2352, Yes base rate 0.378)}\\[1pt]
\mGptFiveTwo & 1{,}480 & 0.2120 & $-$0.0255 {\scriptsize[$-$0.0344, $-$0.0169]} & 0.0044 & 0.0268 & +0.023 {\scriptsize[+0.015, +0.030]} & +0.099 \\
\mGptOss & 1{,}480 & 0.2281 & $-$0.0416 {\scriptsize[$-$0.0517, $-$0.0310]} & 0.0102 & 0.0167 & +0.033 {\scriptsize[+0.024, +0.041]} & +0.030 \\
\mKimiKTwoFive & 1{,}477 & 0.2158 & $-$0.0294 {\scriptsize[$-$0.0389, $-$0.0205]} & 0.0045 & 0.0229 & +0.027 {\scriptsize[+0.018, +0.035]} & +0.082 \\
\mDeepSeekVThreeTwo & 1{,}480 & 0.2138 & $-$0.0272 {\scriptsize[$-$0.0361, $-$0.0182]} & 0.0032 & 0.0249 & +0.025 {\scriptsize[+0.017, +0.032]} & +0.091 \\
\mGptFiveFive & 1{,}480 & 0.2080 & $-$0.0214 {\scriptsize[$-$0.0301, $-$0.0128]} & 0.0036 & 0.0306 & +0.019 {\scriptsize[+0.011, +0.027]} & +0.116 \\
\mGlmFiveOne & 1{,}480 & 0.2149 & $-$0.0284 {\scriptsize[$-$0.0379, $-$0.0189]} & 0.0052 & 0.0253 & +0.024 {\scriptsize[+0.016, +0.033]} & +0.086 \\
\mKimiKTwoSix & 1{,}479 & 0.2126 & $-$0.0264 {\scriptsize[$-$0.0351, $-$0.0180]} & 0.0042 & 0.0262 & +0.023 {\scriptsize[+0.016, +0.031]} & +0.096 \\
\mDeepSeekVFourPro & 1{,}475 & 0.2160 & $-$0.0298 {\scriptsize[$-$0.0401, $-$0.0200]} & 0.0071 & 0.0259 & +0.024 {\scriptsize[+0.016, +0.032]} & +0.082 \\
\addlinespace
\market & --- & 0.1864 & --- & 0.0017 & 0.0500 & --- & +0.207 \\
Market, shrunk & --- & 0.1956 & --- & --- & --- & --- & +0.168 \\
Base rate every row & --- & 0.2352 & --- & --- & 0.0000 & --- & +0.000 \\
0.5 every row & --- & 0.2500 & --- & --- & 0.0000 & --- & $-$0.063 \\
Yes every row & --- & 0.6217 & --- & --- & 0.0000 & --- & $-$1.643 \\
No every row & --- & 0.3783 & --- & --- & 0.0000 & --- & $-$0.608 \\
\addlinespace
\multicolumn{8}{@{}l}{\textbf{\winC} (uncertainty 0.2402, Yes base rate 0.401)}\\[1pt]
\mGptFiveTwo & 2{,}137 & 0.2218 & $-$0.0372 {\scriptsize[$-$0.0447, $-$0.0293]} & 0.0063 & 0.0237 & +0.032 {\scriptsize[+0.025, +0.038]} & +0.077 \\
\mGptOss & 2{,}137 & 0.2278 & $-$0.0431 {\scriptsize[$-$0.0515, $-$0.0347]} & 0.0073 & 0.0194 & +0.036 {\scriptsize[+0.029, +0.043]} & +0.052 \\
\mKimiKTwoFive & 2{,}134 & 0.2205 & $-$0.0361 {\scriptsize[$-$0.0441, $-$0.0285]} & 0.0046 & 0.0244 & +0.031 {\scriptsize[+0.024, +0.039]} & +0.082 \\
\mDeepSeekVThreeTwo & 2{,}137 & 0.2248 & $-$0.0402 {\scriptsize[$-$0.0480, $-$0.0325]} & 0.0055 & 0.0206 & +0.035 {\scriptsize[+0.028, +0.041]} & +0.064 \\
\mGptFiveFive & 2{,}136 & 0.2117 & $-$0.0272 {\scriptsize[$-$0.0341, $-$0.0200]} & 0.0023 & 0.0302 & +0.026 {\scriptsize[+0.019, +0.032]} & +0.119 \\
\mGlmFiveOne & 2{,}137 & 0.2255 & $-$0.0409 {\scriptsize[$-$0.0489, $-$0.0330]} & 0.0066 & 0.0222 & +0.034 {\scriptsize[+0.027, +0.039]} & +0.061 \\
\mKimiKTwoSix & 2{,}136 & 0.2245 & $-$0.0397 {\scriptsize[$-$0.0477, $-$0.0320]} & 0.0058 & 0.0216 & +0.034 {\scriptsize[+0.027, +0.041]} & +0.066 \\
\mDeepSeekVFourPro & 2{,}130 & 0.2264 & $-$0.0418 {\scriptsize[$-$0.0506, $-$0.0334]} & 0.0106 & 0.0244 & +0.031 {\scriptsize[+0.024, +0.038]} & +0.058 \\
\addlinespace
\market & --- & 0.1846 & --- & 0.0006 & 0.0558 & --- & +0.232 \\
Market, shrunk & --- & 0.1936 & --- & --- & --- & --- & +0.194 \\
Base rate every row & --- & 0.2402 & --- & --- & 0.0000 & --- & +0.000 \\
0.5 every row & --- & 0.2500 & --- & --- & 0.0000 & --- & $-$0.041 \\
Yes every row & --- & 0.5990 & --- & --- & 0.0000 & --- & $-$1.494 \\
No every row & --- & 0.4010 & --- & --- & 0.0000 & --- & $-$0.669 \\
\bottomrule
\end{tabular}
\endgroup

\end{table*}

\section{Which file produces which result}
\label{app:repro}
Every number in this paper comes from a released file. Table~\ref{tab:repro}
maps each float and each named analysis to the script or results file that
produces it. Scripts (a name with no extension, or \code{.py}) and
\code{models.json} are in the code repository, each exploratory script with its
saved output; result tables (\code{.csv}) and the other \code{.json} inputs are
in the main dataset. \winC numbers also need the \kalshi market data, which the
dataset card explains how to obtain and \code{build\_db.py} merges.

\begin{table*}[!tbp]
\centering
\caption{Where each result comes from.}
\label{tab:repro}
\begingroup\footnotesize\setlength{\tabcolsep}{4pt}%
\begin{tabular}{@{}p{0.34\linewidth}p{0.62\linewidth}@{}}
\toprule
Result & Produced by \\
\midrule
Table~\ref{tab:benchmark} & \code{benchmark-primary-set-counts}, \code{depth\_cuts.json}, \code{reporting\_line.csv}, \code{expB\_main.csv} \\
Table~\ref{tab:crowd} & \code{direction-and-money-against-crowd-carries-the-loss}, \code{direction-and-money-against-means-buying-the-longshot} \\
Table~\ref{tab:absolute} & \code{baselines-absolute-and-horizon}, \code{calibration-shape-resolution-gap}, \code{baselines-whole-window-edge} \\
Table~\ref{tab:roster} & \code{models.json}, \code{reporting\_line.csv} \\
Figure~\ref{fig:forecast-vs-price} and the Yes-side cost (Section~\ref{sec:results-crowd}) & \code{news-effect-high-price-gap}, \code{calibration-shape-yes-no-asymmetry} (test-split output) \\
Figure~\ref{fig:brier-gap} & \code{expA\_buckets.csv} \\
Figure~\ref{fig:sorting} & \code{calibration-shape-resolution-gap} \\
Figure~\ref{fig:recalibration} & \code{recalibration-crossfit-and-isotonic} \\
Figure~\ref{fig:shrinkage-null} & \code{correlated-errors-shrinkage-null} \\
Figure~\ref{fig:correlated} & \code{correlated-errors-pair-ranking}, \code{expC\_overall.csv} \\
Shared inputs control (Section~\ref{sec:results-correlated}) & \code{correlated-errors-covariate-adjusted} \\
Agreement with and without news (Appendix~\ref{app:newsagree}) & \code{correlated-errors-with-and-without-news}, \code{correlated-errors-news-covariance-or-variance} \\
Depth in accuracy terms (Section~\ref{sec:results-depth}) & \code{expA\_gap\_trend.csv}, \code{expA\_pooled\_verdicts.csv} \\
Learned combination (Section~\ref{sec:results-correlated}) & \code{ensembles-learned-stacking}, \code{ensembles-and-extremizing} \\
Bin-free decomposition (Section~\ref{sec:results-sorting}) & \code{decomposition-binning-bias} \\
Failure imputation and re-ask rate (Section~\ref{sec:models}) & \code{failures-imputation-sensitivity}, \code{probe-news-value-and-reask-rate} \\
Forecast horizon (Appendix~\ref{app:horizon}) & \code{baselines-absolute-and-horizon} \\
Fee and stake sensitivity (Appendix~\ref{app:feesensitivity}) & \code{fees-stake-sensitivity} \\
News placebo (Section~\ref{sec:results-crowd}) & \code{news-toward-away-placebo} \\
Memory probe and value of news (Appendix~\ref{app:memory}, Section~\ref{sec:results-crowd}) & \code{quarantine\_<run>\_summary.json}, \code{window\_c\_quarantine\_counts.csv}, \code{first\_runs\_report.py}, \code{probe-news-value-and-reask-rate} \\
Skill scores and shrunk market (Section~\ref{sec:results-market}) & \code{baselines-absolute-and-horizon} \\
Non-ladder in-play rows and ladders (Section~\ref{sec:results-market}) & \code{secondary\_full\_set.csv}, \code{secondary\_ladders.csv} \\
Wrong together (Appendix~\ref{app:wrongtogether}) & \code{correlated-errors-shrinkage-null}, \code{correlated-errors-shrinkage-null-news-rows}, \code{correlated-errors-covariate-adjusted}, \code{expC\_overall.csv} \\
Newest model (Appendix~\ref{app:newest}) & \code{pair\_gpt-5.2\_vs\_gpt-5.5\_paired.csv}, \code{calibration-shape-model-spread}, \code{decomposition-binning-bias} \\
Same side or opposite (Appendix~\ref{app:againstdetail}) & \code{direction-same-side-share}, \code{direction-and-money-against-means-buying-the-longshot} \\
Events and clusters (Appendix~\ref{app:events}) & \code{event-definition-and-cluster-counts} \\
Split leakage (Appendix~\ref{app:splitcheck}) & \code{split-leakage-event-grouped} \\
Small reliability term (Appendix~\ref{app:relsmall}) & \code{decomposition-why-reliability-is-small} \\
Pairwise correlations (Appendix~\ref{app:pairdetail}) & \code{correlated-errors-pair-ranking} \\
Non-sports gap (Appendix~\ref{app:nonsports}) & \code{topic-and-sports-brier-gap} (test-split output), \code{topic-and-sports-divergence}, \code{calibration-shape-sports-split} \\
Thinking length (Appendix~\ref{app:thinking}) & \code{thinking-and-length-more-thinking-more-disagreement}, \code{thinking-and-length-budget-cap-harmless} \\
Calibration gate (Appendix~\ref{app:gate}) & \code{calibration\_gate.json}, \code{gate\_used.json}, \code{expB\_main.csv} \\
Convergence (Appendix~\ref{app:convergence}) & \code{expB\_main.csv}, \code{convergence-toward-outcome}, \code{convergence-settled-in-window}, \code{convergence-baseline-choice}, \code{convergence-open-market-move} \\
Exploratory scripts (Appendix~\ref{app:exploratory}) & every script in \code{analysis/exploratory/}, each with its saved output in \code{analysis/exploratory/outputs/} \\
Prompt (Appendix~\ref{app:prompt}) & \code{prompt.py} \\
\bottomrule
\end{tabular}
\endgroup

\end{table*}

\section{Fee formulas}
\label{app:feeformula}
On \polymarket the taker fee is a topic-dependent rate times $p(1-p)$: sports
0.05 from July 2026, crypto 0.07, politics and tech 0.04, geopolitics free. On
\kalshi the fee is 0.07 times a per-market multiplier times $P(1-P)$, rounded
up to the next cent per order.

\section{News attachment rules}
\label{app:newsattach}
An article may attach to a row only if its crawl time is strictly before the
forecast moment and no more than fourteen days before it; it must match the
market by name phrase and clear a reranker score of 0.3; at most the three
best-scoring articles attach. Two dateline rules then apply. The benchmark drops
any row whose news carries a dateline two or more days after the news's own
capture time. At call time the harness separately blocks any row whose news
carries a dateline two or more days after the forecast moment, unless that date
was hand-labelled as a scheduled future event rather than a publication date:
22 rows in all (17 on \winA, 3 on \winB, 2 on \winC, at most 0.6 per cent of a
window's in-play rows), labelled from the article text alone, without access to prices, outcomes or model
forecasts.
One day is allowed in both because capture times are UTC and sites print local
dates.

\section{Skill score}
\label{app:skill}
The skill score against a forecaster who answers the window's base rate
$\bar{y}$ on every question is
\begin{equation}
S(p) \;=\; 1 - \frac{\mathrm{BS}(p)}{\mathrm{BS}(\bar{y})},
\label{eq:skill}
\end{equation}
where zero means no better than answering the base rate every time and one means
perfect.

\section{Correlated-error method}
\label{app:correlmethod}
For model $m$ we fit
\begin{equation}
p_i^{(m)} \;=\; a_m + b_m\,q_i + \varepsilon_i^{(m)}
\label{eq:shrinkage}
\end{equation}
by least squares on the probability scale. We call a slope $b_m<1$ shrinkage:
the model moves less than the price does, whether because it hedges or because it
knows less, and its forecasts sit toward the fixed point $a_m/(1-b_m)$. The quantity we report is the mean pairwise
correlation of the fitted residuals,
\begin{equation}
\bar{\rho} \;=\; \binom{M}{2}^{-1} \sum_{m < m'}
  \mathrm{corr}\!\left(\varepsilon^{(m)},\, \varepsilon^{(m')}\right),
\label{eq:residcorr}
\end{equation}
which is agreement that each model's own fitted shrinkage cannot explain.

Two simulated controls appear in Figure~\ref{fig:shrinkage-null}. Both build
forecasters carrying each model's fitted $a_m$ and $b_m$; they differ in where
the noise comes from. The first draws it from a Gaussian at the model's own
fitted residual spread. The second discards the distributional assumption and
resamples each model's \emph{observed} residuals without replacement within a
price band, so the simulated noise has the real shape and the real spread and
only its pairing across models is destroyed. The two agree closely, which is why
the text quotes them as a range. Neither is a demanding test: independent noise
is uncorrelated by construction, so these controls show only that shrinkage
alone cannot manufacture the agreement, not that nothing else could.

The stronger control replaces $b_m q_i$ with a fitted surface that includes the
row's observable characteristics (a sports indicator, topic tags, horizon in
log days and as band indicators, a news indicator and price-band
indicators) and, on rows belonging to events with more than one market, a
separate intercept per event. Section~\ref{sec:results-correlated} reports
$\bar{\rho}$ under each fit.

The raw error correlation, $\mathrm{corr}(p-y)$, is 0.9625 on \winB and 0.9599
on \winC on the test split. We report it because earlier work uses
it~\citep{spiro2026}, but it should not be compared with $\bar{\rho}$: a raw
error correlation rises whenever forecasts move little relative to outcomes.

\section{Wrong together: per-window detail}
\label{app:wrongtogether}
A pair lands on the same incorrect side of even money 26.4 per cent of the time
on \winB against 10.7 per cent expected under independence, a ratio of 2.47, and
26.9 against 13.3 on \winC, a ratio of 2.03. On news-bearing rows the ratios are
2.33 and 2.56 (431 and 594 rows), read against a shrinkage null of 1.67 and 1.78
on those same rows; restricted further to the news-bearing rows the market got
right, they are 3.11 on \winB and 4.10 on \winC against nulls of 1.4 to 1.8. Simulated forecasters sharing only shrinkage already miss together
about 1.9 times as often as independence on all rows of \winA and \winB and about 1.5 times on
\winC, so the margin is narrower than a comparison against independence alone
would suggest. Between 58 and 74 per cent of shared mistakes fall where the
market was also wrong, leaving a quarter to two fifths that are the models' own.

Richer nulls do not close the gap. On rows the market got right, simulated
forecasters that share each model's fit to the price and the row's observable
features miss together 2.3 times the independent rate on \winA, 1.9 to 2.0 on
\winB and 1.5 on \winC, against real ratios of 4.3, 4.0 and 2.7. Adding a
separate intercept per event, on the rows of events with more than one market,
raises the null to 3.5 to 3.6, 2.9 and 2.2 to 2.3, against real ratios on those
rows of 5.9, 4.3 and 3.2. Each null as a share of the real ratio on its own rows
is 55, 47 to 48 and 58 per cent for features (all market-right rows) and 60 to
61, 67 and 69 to 70 per cent with the event (multi-market rows), so the event
appears to absorb part of the excess, though the two shares rest on different
rows. The real ratio still lies
above every simulated range.

\section{Agreement with and without news}
\label{app:newsagree}
Every news-bearing row was asked twice, with the news and with it removed, so
the models' agreement can be measured on the same rows under both. The residual
correlation after each model's fitted line is 0.756 with news against 0.803
without on \winA, 0.707 against 0.763 on \winB and 0.711 against 0.747 on
\winC; the falls of 0.047, 0.055 and 0.036 exclude zero on \winA and \winB and
reach it on \winC (event-resampled, both passes resampled together, on the
1{,}141, 431 and 593 rows every model answered in both passes; on one \winC row
a model has no answer without the news, hence 593 here against 594 in
Appendix~\ref{app:wrongtogether}). The wrong-together ratio on rows the market got right shows no
clear change: $+0.14$, $-0.26$ and $+0.20$, every interval spanning zero. News
therefore removes a slice of the agreement in departures from the price and none
detectable in shared wrong calls. The test bounds only the information our news
carried.

A correlation can fall because the models share less or because each becomes
noisier. Splitting the change settles which (same rows, same resampling). On
\winA the residual standard deviation rises from 0.154 to 0.168 with news, clear
of zero, while the mean pairwise covariance rises slightly: holding the spreads
at their no-news values, the with-news correlation would be 0.90, not 0.76. The
fall on \winA is therefore model-specific noise added by the news, on the one
window whose news was retrieved at run time. On \winB and \winC the spreads barely move
($-0.002$) and the covariance falls, by 0.0026 and 0.0020, which alone would
lower the correlation by 0.07 and 0.06; neither covariance change clears zero.
The shared-missing-information reading is therefore suggestive on the two
live-book windows and unsupported on \winA.

\section{The two newest models}
\label{app:newest}
\mGptFiveFive has the highest resolution in the study, 0.0306 on \winB and
0.0302 on \winC, and is still short of the market by $+0.019$ and $+0.026$ with
both intervals clear of zero. Against \mGptFiveTwo it is better on \kalshi by
0.0099 (interval excluding zero) and not clearly better on \polymarket.
\mDeepSeekVFourPro shares the same release date yet on \winC has a resolution gap
to the market of $+0.0314$ against \mGptFiveFive's $+0.0257$, so recency does
not separate models released on the same day. \mDeepSeekVFourPro is a capped
model whose host truncated some outputs, so one pair is thin evidence either way.

\section{Leaning against the crowd: where the loss sits}
\label{app:againstdetail}
Most rows that lean against the crowd are same-side-less-confident rather than
outright contradiction, but the accuracy loss is not spread evenly across the
two. On the rows where the model genuinely favours the other outcome the Brier
edge is $-0.084$ on \winA, $-0.121$ on \winB and $-0.104$ on \winC, against
$-0.016$, $-0.013$ and $-0.019$ on the same-side rows: five to ten times worse.
The direction split therefore tracks something real, even though the majority of
the category is a difference of confidence. In money the contrast is the same
shape: trades placed against the crowd buy at 0.28 to 0.31 and win 27 to 29 per
cent of the time, while trades placed with the crowd pay 0.70 to 0.71 and win 70
to 72 per cent. These averages look close to fair, but returns average each
trade's own payoff, and on \winA, where fees were near zero, against-the-crowd
trades still lost 17 cents per dollar.

In the lowest price band the models answer 0.20 to 0.25 against the crowd's 0.13
to 0.15. Each model's fixed point sits near the window's base rate: 0.26 to 0.34
on \winA (base rate 0.37), 0.32 to 0.38 on \winB (0.38) and 0.38 to 0.44 on
\winC (0.40).

\section{What an event is, and how many clusters}
\label{app:events}
Every interval in this paper resamples \emph{events} rather than rows. An event
is the venue's own grouping of markets about one happening (a \polymarket event
id, a \kalshi event ticker), so one event may carry nineteen markets on which
songs are played at a single halftime show, or twelve on what one official says
at one press conference. Markets inside an event move together, so treating them
as independent draws would give intervals that are too narrow.

On the test split the bootstrap draws from 626 events on \winA, 1{,}196 on \winB
and 1{,}823 on \winC, holding 1{,}248, 1{,}480 and 2{,}137 market-moments. Markets
per event in the test split run from one to thirteen. How much the clustering
matters depends on the window: 46.6 per cent of \winA events carry a single market,
against 83.5 per cent on \winB and 89.1 per cent on \winC. On the two larger
windows the event bootstrap is therefore close to a row bootstrap, and the
protection it offers is real but small; on \winA it does substantial work.

\section{Does the calibration split leak across events?}
\label{app:splitcheck}
The calibration and test split was drawn per market rather than per event, so
two markets from one event can land on opposite sides. Since the recalibration
control fits a map on one side and scores it on the other, a map might be fitted
on rows closely related to those it is scored on. The exposure is real:
straddling events hold 31.3 per cent of \winA test rows, 15.7 per cent of \winB
and 10.2 per cent of \winC.

We therefore redrew the split, same cells and same size, keeping every event
whole, and refitted the control. Recovery does not rise materially. On \winB it runs
from $-0.5$ to 17.1 per cent against 1.7 to 16.1 per cent under the main split; on \winC every model recovers \emph{less}, the range falling from 5.7 to
17.9 per cent down to 0.6 to 9.9 per cent. Scoring the main-split map on
non-straddling rows alone moves nothing systematically. The cross-fitted Platt
and isotonic figures, which set the paper's ceiling of about a fifth, are
unaffected either way because those folds are built by event. The
conclusion that no monotone rescaling closes the gap is unchanged, and on \winC
it is stronger.

\section{Why the reliability term is small}
\label{app:relsmall}
Reliability being small does not mean the models are well calibrated, and three
separate things keep it small. First, it is a squared quantity: a per-model
reliability of 0.0023 to 0.0106 is a root-mean-square miss of roughly 0.05 to
0.10 in probability, set against a market-minus-model resolution gap of 0.019 to 0.036.
Second, the badly calibrated forecasts sit in the thin tails. About three in ten
forecasts fall outside 0.2 to 0.7, and those carry 68 per cent of the total
reliability on \winA, 95 per cent on \winB and 78 per cent on \winC, while the
crowded middle bins miss by only 0.012 to 0.039. Third, ten equal-width bins
hide disagreement inside a bin. Splitting each bin by whether the market priced
the row above or below the bin's mean forecast multiplies the pooled
reliability by three to five, from 0.0027 to 0.0128 on \winA, 0.0043 to 0.0130 on \winB and
0.0055 to 0.0191 on \winC; a random split into halves of the same size leaves it
almost unchanged. That third figure is a diagnostic, not a competing
decomposition: Murphy's reliability conditions on the forecast alone, so
conditioning on the price as well moves part of the sorting shortfall into the
calibration column by construction. Measured against a model's own information
the shortfall is a sorting gap; measured against the market's better information
the same fact appears as miscalibration. The 79 to 99 per cent attribution in
Section~\ref{sec:results-sorting} is the former. Note that the market sits in the 0.2 to 0.7 range about as
often as the models do, with reliability three to nine times smaller, so the
range alone explains nothing.

\section{Fee and stake sensitivity}
\label{app:feesensitivity}
The same trades with no fee at all return $-0.008$ on \winB against $-0.033$
with fees, and $-0.059$ on \winC against $-0.099$: before costs the \winB trades
are roughly break-even. On \winA fees were about 0.02 per cent of outlay, so the
return of $-0.102$ is the same with or without them. Raising the stake from \$1 to \$1{,}000 moves the
\kalshi return by 0.4 percentage points, so the result is not an artefact of
small orders.

\section{Pairwise departure correlations}
\label{app:pairdetail}
\begin{figure*}[!tbp]
\centering
\begin{tikzpicture}
\begin{groupplot}[
  paperaxis,
  group style={group size=2 by 1, horizontal sep=6pt},
  width=0.40\linewidth, height=0.40\linewidth,
  xmin=0, xmax=8, ymin=0, ymax=8,
  axis equal image,
  xtick={0.5,1.5,2.5,3.5,4.5,5.5,6.5,7.5}, ytick={0.5,1.5,2.5,3.5,4.5,5.5,6.5,7.5},
  xticklabels={{\mGptFiveTwo},{\mGptOss},{\mKimiKTwoFive},{\mDeepSeekVThreeTwo},{\mGptFiveFive},{\mGlmFiveOne},{\mKimiKTwoSix},{\mDeepSeekVFourPro}},
  x tick label style={rotate=60, anchor=east, font=\scriptsize},
  y tick label style={font=\scriptsize},
  tick style={draw=none},
]
\nextgroupplot[title={\winB (\polymarket)}, yticklabels={{\mDeepSeekVFourPro},{\mKimiKTwoSix},{\mGlmFiveOne},{\mGptFiveFive},{\mDeepSeekVThreeTwo},{\mKimiKTwoFive},{\mGptOss},{\mGptFiveTwo}},]
\fill[fill=black!42] (axis cs:1,7) rectangle (axis cs:2,8);
\node[font=\tiny, text=black] at (axis cs:1.5,7.5) {.75};
\fill[fill=black!75] (axis cs:2,7) rectangle (axis cs:3,8);
\node[font=\tiny, text=white] at (axis cs:2.5,7.5) {.86};
\fill[fill=black!67] (axis cs:3,7) rectangle (axis cs:4,8);
\node[font=\tiny, text=white] at (axis cs:3.5,7.5) {.83};
\fill[fill=black!78] (axis cs:4,7) rectangle (axis cs:5,8);
\node[font=\tiny, text=white] at (axis cs:4.5,7.5) {.87};
\fill[fill=black!65] (axis cs:5,7) rectangle (axis cs:6,8);
\node[font=\tiny, text=white] at (axis cs:5.5,7.5) {.83};
\fill[fill=black!73] (axis cs:6,7) rectangle (axis cs:7,8);
\node[font=\tiny, text=white] at (axis cs:6.5,7.5) {.85};
\fill[fill=black!56] (axis cs:7,7) rectangle (axis cs:8,8);
\node[font=\tiny, text=black] at (axis cs:7.5,7.5) {.80};
\fill[fill=black!42] (axis cs:0,6) rectangle (axis cs:1,7);
\node[font=\tiny, text=black] at (axis cs:0.5,6.5) {.75};
\fill[fill=black!44] (axis cs:2,6) rectangle (axis cs:3,7);
\node[font=\tiny, text=black] at (axis cs:2.5,6.5) {.75};
\fill[fill=black!42] (axis cs:3,6) rectangle (axis cs:4,7);
\node[font=\tiny, text=black] at (axis cs:3.5,6.5) {.75};
\fill[fill=black!19] (axis cs:4,6) rectangle (axis cs:5,7);
\node[font=\tiny, text=black] at (axis cs:4.5,6.5) {.67};
\fill[fill=black!30] (axis cs:5,6) rectangle (axis cs:6,7);
\node[font=\tiny, text=black] at (axis cs:5.5,6.5) {.71};
\fill[fill=black!45] (axis cs:6,6) rectangle (axis cs:7,7);
\node[font=\tiny, text=black] at (axis cs:6.5,6.5) {.76};
\fill[fill=black!16] (axis cs:7,6) rectangle (axis cs:8,7);
\node[font=\tiny, text=black] at (axis cs:7.5,6.5) {.66};
\fill[fill=black!75] (axis cs:0,5) rectangle (axis cs:1,6);
\node[font=\tiny, text=white] at (axis cs:0.5,5.5) {.86};
\fill[fill=black!44] (axis cs:1,5) rectangle (axis cs:2,6);
\node[font=\tiny, text=black] at (axis cs:1.5,5.5) {.75};
\fill[fill=black!69] (axis cs:3,5) rectangle (axis cs:4,6);
\node[font=\tiny, text=white] at (axis cs:3.5,5.5) {.84};
\fill[fill=black!59] (axis cs:4,5) rectangle (axis cs:5,6);
\node[font=\tiny, text=white] at (axis cs:4.5,5.5) {.81};
\fill[fill=black!67] (axis cs:5,5) rectangle (axis cs:6,6);
\node[font=\tiny, text=white] at (axis cs:5.5,5.5) {.83};
\fill[fill=black!89] (axis cs:6,5) rectangle (axis cs:7,6);
\node[font=\tiny, text=white] at (axis cs:6.5,5.5) {.91};
\fill[fill=black!59] (axis cs:7,5) rectangle (axis cs:8,6);
\node[font=\tiny, text=white] at (axis cs:7.5,5.5) {.81};
\fill[fill=black!67] (axis cs:0,4) rectangle (axis cs:1,5);
\node[font=\tiny, text=white] at (axis cs:0.5,4.5) {.83};
\fill[fill=black!42] (axis cs:1,4) rectangle (axis cs:2,5);
\node[font=\tiny, text=black] at (axis cs:1.5,4.5) {.75};
\fill[fill=black!69] (axis cs:2,4) rectangle (axis cs:3,5);
\node[font=\tiny, text=white] at (axis cs:2.5,4.5) {.84};
\fill[fill=black!50] (axis cs:4,4) rectangle (axis cs:5,5);
\node[font=\tiny, text=black] at (axis cs:4.5,4.5) {.78};
\fill[fill=black!55] (axis cs:5,4) rectangle (axis cs:6,5);
\node[font=\tiny, text=black] at (axis cs:5.5,4.5) {.79};
\fill[fill=black!68] (axis cs:6,4) rectangle (axis cs:7,5);
\node[font=\tiny, text=white] at (axis cs:6.5,4.5) {.84};
\fill[fill=black!45] (axis cs:7,4) rectangle (axis cs:8,5);
\node[font=\tiny, text=black] at (axis cs:7.5,4.5) {.76};
\fill[fill=black!78] (axis cs:0,3) rectangle (axis cs:1,4);
\node[font=\tiny, text=white] at (axis cs:0.5,3.5) {.87};
\fill[fill=black!19] (axis cs:1,3) rectangle (axis cs:2,4);
\node[font=\tiny, text=black] at (axis cs:1.5,3.5) {.67};
\fill[fill=black!59] (axis cs:2,3) rectangle (axis cs:3,4);
\node[font=\tiny, text=white] at (axis cs:2.5,3.5) {.81};
\fill[fill=black!50] (axis cs:3,3) rectangle (axis cs:4,4);
\node[font=\tiny, text=black] at (axis cs:3.5,3.5) {.78};
\fill[fill=black!54] (axis cs:5,3) rectangle (axis cs:6,4);
\node[font=\tiny, text=black] at (axis cs:5.5,3.5) {.79};
\fill[fill=black!63] (axis cs:6,3) rectangle (axis cs:7,4);
\node[font=\tiny, text=white] at (axis cs:6.5,3.5) {.82};
\fill[fill=black!50] (axis cs:7,3) rectangle (axis cs:8,4);
\node[font=\tiny, text=black] at (axis cs:7.5,3.5) {.77};
\fill[fill=black!65] (axis cs:0,2) rectangle (axis cs:1,3);
\node[font=\tiny, text=white] at (axis cs:0.5,2.5) {.83};
\fill[fill=black!30] (axis cs:1,2) rectangle (axis cs:2,3);
\node[font=\tiny, text=black] at (axis cs:1.5,2.5) {.71};
\fill[fill=black!67] (axis cs:2,2) rectangle (axis cs:3,3);
\node[font=\tiny, text=white] at (axis cs:2.5,2.5) {.83};
\fill[fill=black!55] (axis cs:3,2) rectangle (axis cs:4,3);
\node[font=\tiny, text=black] at (axis cs:3.5,2.5) {.79};
\fill[fill=black!54] (axis cs:4,2) rectangle (axis cs:5,3);
\node[font=\tiny, text=black] at (axis cs:4.5,2.5) {.79};
\fill[fill=black!65] (axis cs:6,2) rectangle (axis cs:7,3);
\node[font=\tiny, text=white] at (axis cs:6.5,2.5) {.83};
\fill[fill=black!57] (axis cs:7,2) rectangle (axis cs:8,3);
\node[font=\tiny, text=black] at (axis cs:7.5,2.5) {.80};
\fill[fill=black!73] (axis cs:0,1) rectangle (axis cs:1,2);
\node[font=\tiny, text=white] at (axis cs:0.5,1.5) {.85};
\fill[fill=black!45] (axis cs:1,1) rectangle (axis cs:2,2);
\node[font=\tiny, text=black] at (axis cs:1.5,1.5) {.76};
\fill[fill=black!89] (axis cs:2,1) rectangle (axis cs:3,2);
\node[font=\tiny, text=white] at (axis cs:2.5,1.5) {.91};
\fill[fill=black!68] (axis cs:3,1) rectangle (axis cs:4,2);
\node[font=\tiny, text=white] at (axis cs:3.5,1.5) {.84};
\fill[fill=black!63] (axis cs:4,1) rectangle (axis cs:5,2);
\node[font=\tiny, text=white] at (axis cs:4.5,1.5) {.82};
\fill[fill=black!65] (axis cs:5,1) rectangle (axis cs:6,2);
\node[font=\tiny, text=white] at (axis cs:5.5,1.5) {.83};
\fill[fill=black!59] (axis cs:7,1) rectangle (axis cs:8,2);
\node[font=\tiny, text=white] at (axis cs:7.5,1.5) {.81};
\fill[fill=black!56] (axis cs:0,0) rectangle (axis cs:1,1);
\node[font=\tiny, text=black] at (axis cs:0.5,0.5) {.80};
\fill[fill=black!16] (axis cs:1,0) rectangle (axis cs:2,1);
\node[font=\tiny, text=black] at (axis cs:1.5,0.5) {.66};
\fill[fill=black!59] (axis cs:2,0) rectangle (axis cs:3,1);
\node[font=\tiny, text=white] at (axis cs:2.5,0.5) {.81};
\fill[fill=black!45] (axis cs:3,0) rectangle (axis cs:4,1);
\node[font=\tiny, text=black] at (axis cs:3.5,0.5) {.76};
\fill[fill=black!50] (axis cs:4,0) rectangle (axis cs:5,1);
\node[font=\tiny, text=black] at (axis cs:4.5,0.5) {.77};
\fill[fill=black!57] (axis cs:5,0) rectangle (axis cs:6,1);
\node[font=\tiny, text=black] at (axis cs:5.5,0.5) {.80};
\fill[fill=black!59] (axis cs:6,0) rectangle (axis cs:7,1);
\node[font=\tiny, text=white] at (axis cs:6.5,0.5) {.81};
\draw[gray!40] (axis cs:0,7) -- (axis cs:1,8);
\draw[gray!40] (axis cs:1,6) -- (axis cs:2,7);
\draw[gray!40] (axis cs:2,5) -- (axis cs:3,6);
\draw[gray!40] (axis cs:3,4) -- (axis cs:4,5);
\draw[gray!40] (axis cs:4,3) -- (axis cs:5,4);
\draw[gray!40] (axis cs:5,2) -- (axis cs:6,3);
\draw[gray!40] (axis cs:6,1) -- (axis cs:7,2);
\draw[gray!40] (axis cs:7,0) -- (axis cs:8,1);
\nextgroupplot[title={\winC (\kalshi)}, yticklabels={},]
\fill[fill=black!25] (axis cs:1,7) rectangle (axis cs:2,8);
\node[font=\tiny, text=black] at (axis cs:1.5,7.5) {.69};
\fill[fill=black!59] (axis cs:2,7) rectangle (axis cs:3,8);
\node[font=\tiny, text=white] at (axis cs:2.5,7.5) {.81};
\fill[fill=black!55] (axis cs:3,7) rectangle (axis cs:4,8);
\node[font=\tiny, text=black] at (axis cs:3.5,7.5) {.79};
\fill[fill=black!57] (axis cs:4,7) rectangle (axis cs:5,8);
\node[font=\tiny, text=black] at (axis cs:4.5,7.5) {.80};
\fill[fill=black!64] (axis cs:5,7) rectangle (axis cs:6,8);
\node[font=\tiny, text=white] at (axis cs:5.5,7.5) {.82};
\fill[fill=black!57] (axis cs:6,7) rectangle (axis cs:7,8);
\node[font=\tiny, text=white] at (axis cs:6.5,7.5) {.80};
\fill[fill=black!54] (axis cs:7,7) rectangle (axis cs:8,8);
\node[font=\tiny, text=black] at (axis cs:7.5,7.5) {.79};
\fill[fill=black!25] (axis cs:0,6) rectangle (axis cs:1,7);
\node[font=\tiny, text=black] at (axis cs:0.5,6.5) {.69};
\fill[fill=black!42] (axis cs:2,6) rectangle (axis cs:3,7);
\node[font=\tiny, text=black] at (axis cs:2.5,6.5) {.75};
\fill[fill=black!44] (axis cs:3,6) rectangle (axis cs:4,7);
\node[font=\tiny, text=black] at (axis cs:3.5,6.5) {.75};
\fill[fill=black!7] (axis cs:4,6) rectangle (axis cs:5,7);
\node[font=\tiny, text=black] at (axis cs:4.5,6.5) {.63};
\fill[fill=black!32] (axis cs:5,6) rectangle (axis cs:6,7);
\node[font=\tiny, text=black] at (axis cs:5.5,6.5) {.71};
\fill[fill=black!42] (axis cs:6,6) rectangle (axis cs:7,7);
\node[font=\tiny, text=black] at (axis cs:6.5,6.5) {.75};
\fill[fill=black!17] (axis cs:7,6) rectangle (axis cs:8,7);
\node[font=\tiny, text=black] at (axis cs:7.5,6.5) {.66};
\fill[fill=black!59] (axis cs:0,5) rectangle (axis cs:1,6);
\node[font=\tiny, text=white] at (axis cs:0.5,5.5) {.81};
\fill[fill=black!42] (axis cs:1,5) rectangle (axis cs:2,6);
\node[font=\tiny, text=black] at (axis cs:1.5,5.5) {.75};
\fill[fill=black!69] (axis cs:3,5) rectangle (axis cs:4,6);
\node[font=\tiny, text=white] at (axis cs:3.5,5.5) {.84};
\fill[fill=black!47] (axis cs:4,5) rectangle (axis cs:5,6);
\node[font=\tiny, text=black] at (axis cs:4.5,5.5) {.76};
\fill[fill=black!65] (axis cs:5,5) rectangle (axis cs:6,6);
\node[font=\tiny, text=white] at (axis cs:5.5,5.5) {.83};
\fill[fill=black!83] (axis cs:6,5) rectangle (axis cs:7,6);
\node[font=\tiny, text=white] at (axis cs:6.5,5.5) {.89};
\fill[fill=black!59] (axis cs:7,5) rectangle (axis cs:8,6);
\node[font=\tiny, text=white] at (axis cs:7.5,5.5) {.81};
\fill[fill=black!55] (axis cs:0,4) rectangle (axis cs:1,5);
\node[font=\tiny, text=black] at (axis cs:0.5,4.5) {.79};
\fill[fill=black!44] (axis cs:1,4) rectangle (axis cs:2,5);
\node[font=\tiny, text=black] at (axis cs:1.5,4.5) {.75};
\fill[fill=black!69] (axis cs:2,4) rectangle (axis cs:3,5);
\node[font=\tiny, text=white] at (axis cs:2.5,4.5) {.84};
\fill[fill=black!45] (axis cs:4,4) rectangle (axis cs:5,5);
\node[font=\tiny, text=black] at (axis cs:4.5,4.5) {.76};
\fill[fill=black!57] (axis cs:5,4) rectangle (axis cs:6,5);
\node[font=\tiny, text=white] at (axis cs:5.5,4.5) {.80};
\fill[fill=black!70] (axis cs:6,4) rectangle (axis cs:7,5);
\node[font=\tiny, text=white] at (axis cs:6.5,4.5) {.85};
\fill[fill=black!48] (axis cs:7,4) rectangle (axis cs:8,5);
\node[font=\tiny, text=black] at (axis cs:7.5,4.5) {.77};
\fill[fill=black!57] (axis cs:0,3) rectangle (axis cs:1,4);
\node[font=\tiny, text=black] at (axis cs:0.5,3.5) {.80};
\fill[fill=black!7] (axis cs:1,3) rectangle (axis cs:2,4);
\node[font=\tiny, text=black] at (axis cs:1.5,3.5) {.63};
\fill[fill=black!47] (axis cs:2,3) rectangle (axis cs:3,4);
\node[font=\tiny, text=black] at (axis cs:2.5,3.5) {.76};
\fill[fill=black!45] (axis cs:3,3) rectangle (axis cs:4,4);
\node[font=\tiny, text=black] at (axis cs:3.5,3.5) {.76};
\fill[fill=black!42] (axis cs:5,3) rectangle (axis cs:6,4);
\node[font=\tiny, text=black] at (axis cs:5.5,3.5) {.75};
\fill[fill=black!46] (axis cs:6,3) rectangle (axis cs:7,4);
\node[font=\tiny, text=black] at (axis cs:6.5,3.5) {.76};
\fill[fill=black!36] (axis cs:7,3) rectangle (axis cs:8,4);
\node[font=\tiny, text=black] at (axis cs:7.5,3.5) {.73};
\fill[fill=black!64] (axis cs:0,2) rectangle (axis cs:1,3);
\node[font=\tiny, text=white] at (axis cs:0.5,2.5) {.82};
\fill[fill=black!32] (axis cs:1,2) rectangle (axis cs:2,3);
\node[font=\tiny, text=black] at (axis cs:1.5,2.5) {.71};
\fill[fill=black!65] (axis cs:2,2) rectangle (axis cs:3,3);
\node[font=\tiny, text=white] at (axis cs:2.5,2.5) {.83};
\fill[fill=black!57] (axis cs:3,2) rectangle (axis cs:4,3);
\node[font=\tiny, text=white] at (axis cs:3.5,2.5) {.80};
\fill[fill=black!42] (axis cs:4,2) rectangle (axis cs:5,3);
\node[font=\tiny, text=black] at (axis cs:4.5,2.5) {.75};
\fill[fill=black!62] (axis cs:6,2) rectangle (axis cs:7,3);
\node[font=\tiny, text=white] at (axis cs:6.5,2.5) {.82};
\fill[fill=black!61] (axis cs:7,2) rectangle (axis cs:8,3);
\node[font=\tiny, text=white] at (axis cs:7.5,2.5) {.81};
\fill[fill=black!57] (axis cs:0,1) rectangle (axis cs:1,2);
\node[font=\tiny, text=white] at (axis cs:0.5,1.5) {.80};
\fill[fill=black!42] (axis cs:1,1) rectangle (axis cs:2,2);
\node[font=\tiny, text=black] at (axis cs:1.5,1.5) {.75};
\fill[fill=black!83] (axis cs:2,1) rectangle (axis cs:3,2);
\node[font=\tiny, text=white] at (axis cs:2.5,1.5) {.89};
\fill[fill=black!70] (axis cs:3,1) rectangle (axis cs:4,2);
\node[font=\tiny, text=white] at (axis cs:3.5,1.5) {.85};
\fill[fill=black!46] (axis cs:4,1) rectangle (axis cs:5,2);
\node[font=\tiny, text=black] at (axis cs:4.5,1.5) {.76};
\fill[fill=black!62] (axis cs:5,1) rectangle (axis cs:6,2);
\node[font=\tiny, text=white] at (axis cs:5.5,1.5) {.82};
\fill[fill=black!56] (axis cs:7,1) rectangle (axis cs:8,2);
\node[font=\tiny, text=black] at (axis cs:7.5,1.5) {.80};
\fill[fill=black!54] (axis cs:0,0) rectangle (axis cs:1,1);
\node[font=\tiny, text=black] at (axis cs:0.5,0.5) {.79};
\fill[fill=black!17] (axis cs:1,0) rectangle (axis cs:2,1);
\node[font=\tiny, text=black] at (axis cs:1.5,0.5) {.66};
\fill[fill=black!59] (axis cs:2,0) rectangle (axis cs:3,1);
\node[font=\tiny, text=white] at (axis cs:2.5,0.5) {.81};
\fill[fill=black!48] (axis cs:3,0) rectangle (axis cs:4,1);
\node[font=\tiny, text=black] at (axis cs:3.5,0.5) {.77};
\fill[fill=black!36] (axis cs:4,0) rectangle (axis cs:5,1);
\node[font=\tiny, text=black] at (axis cs:4.5,0.5) {.73};
\fill[fill=black!61] (axis cs:5,0) rectangle (axis cs:6,1);
\node[font=\tiny, text=white] at (axis cs:5.5,0.5) {.81};
\fill[fill=black!56] (axis cs:6,0) rectangle (axis cs:7,1);
\node[font=\tiny, text=black] at (axis cs:6.5,0.5) {.80};
\draw[gray!40] (axis cs:0,7) -- (axis cs:1,8);
\draw[gray!40] (axis cs:1,6) -- (axis cs:2,7);
\draw[gray!40] (axis cs:2,5) -- (axis cs:3,6);
\draw[gray!40] (axis cs:3,4) -- (axis cs:4,5);
\draw[gray!40] (axis cs:4,3) -- (axis cs:5,4);
\draw[gray!40] (axis cs:5,2) -- (axis cs:6,3);
\draw[gray!40] (axis cs:6,1) -- (axis cs:7,2);
\draw[gray!40] (axis cs:7,0) -- (axis cs:8,1);
\end{groupplot}
\end{tikzpicture}

\caption{Pairwise correlation of departures from the market price, $p-q$, on the rows all
eight forecast. The mean across pairs is 0.78 to 0.80, against 0.23 to 0.28 for simulated
forecasters sharing only shrinkage; after removing shrinkage the residual correlation is 0.71
to 0.75 against essentially zero, which is zero by construction (Appendix~\ref{app:correlmethod}). Test split of the primary set.}
\label{fig:correlated}
\end{figure*}

Among the five same-developer pairs, two are the tightest in the roster:
\mKimiKTwoFive with \mKimiKTwoSix at 0.910 on \winB and 0.891 on \winC, and
\mGptFiveTwo with \mGptFiveFive at 0.872 on \winB, which is what a
same-developer story predicts. The two DeepSeek models are no closer than
unrelated pairs, sitting 0.061 and 0.026 below the roster reference;
\mGptOss's pairs sit 0.088, 0.099 and 0.090 below it and are six of the seven
loosest on \winB. On two windows of three the loosest pair is itself
same-developer: \mGptFiveTwo with \mGptOss at 0.722 on \winA and \mGptFiveFive
with \mGptOss at 0.626 on \winC.

\section{The non-sports gap}
\label{app:nonsports}
Sports dominate all three windows, so it matters that the model-to-market gap is
not a sports artefact. It is not: on the non-sports minority the gap is several
times larger. Pooling the roster, the Brier edge on sports rows is $-0.0061$, $-0.0131$ and
$-0.0266$ on \winA, \winB and \winC, against $-0.0739$, $-0.0727$ and $-0.0873$
on non-sports rows. The sports-minus-non-sports difference is $+0.068$, $+0.060$
and $+0.061$, with every interval clear of zero, and each model's own difference
clears zero on every window. Per model the ratio of the two gaps runs from 2.4 to
6.7 on \winB and \winC. On \winA the ratios are 4.9, 11.7 and 34.0 and one is
undefined, because the sports gap there is
so near zero (\mKimiKTwoFive's is slightly positive) that the ratio stops being
meaningful; the difference, not the ratio, is the stable quantity on \winA. The practical reading is that the sports
concentration described in Section~\ref{sec:limitations} makes the models look
better than they are, not worse.

Sorting and calibration hold up within each half. Run separately on the test
split's sports and non-sports rows, the bin-free decomposition puts 85 to 109 per
cent of the shortfall on resolution for sports rows on \winB and 87 to 96 per cent
on \winC (above 100 where the model is better calibrated than the market), and
cross-fitted rescaling recovers at most 18 and 11 per cent. On non-sports rows
resolution still carries most of it, 73 to 90 per cent on \winB and 59 to 89 per
cent on \winC, but rescaling recovers more: up to 30 and 42 per cent, largest for
\mDeepSeekVFourPro and, on \winC, \mGlmFiveOne, the worst-calibrated models on
those rows. On \winA the sports gap is so small ($-0.001$ to $+0.018$) that the
shares are unstable; on its non-sports rows resolution carries 86 to 97 per cent
and rescaling recovers at most 20. Calibration therefore matters more where the
gap is widest, but it is not what the gap is made of.

\section{Thinking length within a model}
\label{app:thinking}
Comparing thinking budgets \emph{between} models is confounded
(Section~\ref{sec:limitations}); comparing a model with itself removes that
confound but not the difficulty of the rows.
Splitting each model's rows into its own top and bottom thinking quartile, the
rows on which a model thought longest are those on which it departed further from
the market and scored worse, not better. For \mGptFiveTwo the departure
$|p-q|$ rises from 0.071 to 0.159 between bottom and top quartile on \winA, with
the Brier edge falling from $-0.0018$ to $-0.0230$; the pattern repeats on \winB.
\mGptFiveFive behaves the same way, with departures $+0.085$ (\winB) and $+0.041$
(\winC) higher in its top quartile and Brier edges of $-0.0372$ and $-0.0462$
there, against $-0.0049$ and $-0.0201$ in its bottom quartile. This comparison
uses both splits. Longer deliberation is associated with more disagreement with the crowd,
and in this study disagreement with the crowd is what costs. We read this as
selection rather than causation: hard rows invite more thinking.

\section{The calibration gate}
\label{app:gate}
A disagreement with the price is only informative where a forecaster's
probabilities mean something, so each model-window-bucket cell is admitted to
the convergence test only if it holds at least 30 rows and its expected calibration error, computed over at most
three bins of near-equal count, is at most 0.10; a cell with fewer than 30 rows
stays closed. The gate is applied per cell rather than per
model, so a model may enter the test in the buckets where it is calibrated and not in
others. Failures are not concentrated in one model: in most failing cells the
sports rows pass on their own and the non-sports rows are badly miscalibrated,
the same
asymmetry as Appendix~\ref{app:nonsports}. Per-cell verdicts
and admitted-row counts are in the released results tables.

\section{Convergence: do prices drift toward the model?}
\label{app:convergence}
A separate question is whether, when a model disagrees with the price, the price
later moves toward the model. Measuring the share of disagreements whose
subsequent price path drifts model-ward, and comparing it with the same statistic
after shuffling which model disagreed on which row, the answer depends on the
window and does not replicate. On \winA the drift exceeds the shuffled
baseline for all four models, by $+0.095$ to $+0.111$. On \winB the difference
runs from $-0.039$ to $+0.042$ across the eight models, spanning zero, and on
\winC it runs from $-0.048$ to $-0.001$, that is, slightly the wrong way. We
therefore make no claim that prices move toward these models. The \winA result is
the one that would need explaining, and \winA is the window drawn without the
live-book rule.

\section{Forecast horizon}
\label{app:horizon}
The size of the deficit depends strongly on how far ahead the forecast sits; its
direction does not replicate across windows. Splitting rows by time to
settlement, the pooled Brier edge varies by a factor of two to four within each
window. On \winB the worst band is beyond 7 days ($-0.063$ against $-0.016$ at 3
to 7 days), and this is not an averaging artefact: beyond 7 days is the worst
band for all eight models. On \winC the worst band is under a day ($-0.052$),
again for all eight models. On \winA the worst is 1 to 3 days ($-0.041$). So
horizon matters within a window and the models agree about how, but which end of
the horizon hurts reverses between \winB and \winC, and we cannot say why.

One band is worth naming because it is the exception in the study. On \winA
under a day the pooled edge is $+0.0229$, with an interval that crosses zero;
for one model, \mKimiKTwoFive, it is $+0.0452$ with an interval above zero. That
is the only model-and-band cell anywhere in this work where a model is clearly
more accurate than the price. It rests on 86 rows, on \winA, where only four models are eligible and on which we rest no claim
(Section~\ref{sec:results-market}), so we report it rather than build on it.

\section{Exploratory scripts}
\label{app:exploratory}
Every exploratory script in the release was run over all eight models and is
released with its saved output; the appendix sections above cite the ones they
use.

\section{Roster, hosts, precision, thinking, failures and cost}
\label{app:roster}
Table~\ref{tab:roster} lists the eight scored models with the settings each was
run under. Eligibility is by release date, as Section~\ref{sec:models} describes:
the four models published between August 2025 and January 2026 precede all
three windows, while the four published in April 2026 precede \winC and \winB
only. The validation model used in the memory
probe (Appendix~\ref{app:memory}) is never scored and is not listed. Across the whole study, 97{,}465 model calls were made, at a cost of about \$620.

\begin{table*}[!tbp]
\centering
\caption{The eight scored models. \emph{Windows} lists the windows each model is eligible for. \emph{Thinking} is either a reasoning-effort level (the OpenAI models and
\mGptOss) or a 1{,}000-token thinking budget; median thinking is tokens per reply, over every
recorded reply of the run. The reply limit is 4{,}000 tokens for every model, although
\mDeepSeekVFourPro's host truncated output at 3{,}000 on about 6\% of its calls. \emph{Failed
rows} gave no usable probability after one stricter re-ask; counts are for the scored pass,
and the memory-probe pass failed on a further 3 rows for \mKimiKTwoFive, 7 for
\mDeepSeekVFourPro and 1 for \mKimiKTwoSix. A dash marks a window a model was not eligible
for, or a precision its host does not state. Sources: \code{models.json},
\code{reporting\_line.csv}.}
\label{tab:roster}
\begingroup\footnotesize\setlength{\tabcolsep}{3pt}%
\begin{tabular}{@{}llllllll@{}}
\toprule
Model & Windows & Released & Host & Precision & Thinking & Median thinking A / B / C & Failed rows (all in play) A / B / C \\
\midrule
\mGptFiveTwo & A, B, C & 2025-12-10 & OpenAI & -- & effort medium & 233 / 198 / 223 & 0 / 0 / 0 \\
\mGptOss & A, B, C & 2025-08-05 & AkashML & bf16 & effort medium & 264 / 218 / 219 & 0 / 0 / 0 \\
\mKimiKTwoFive & A, B, C & 2026-01-27 & SiliconFlow & int4 & cap 1{,}000 tokens & 1{,}000 / 1{,}000 / 1{,}000 & 28 / 7 / 6 \\
\mDeepSeekVThreeTwo & A, B, C & 2025-12-01 & SiliconFlow & fp8 & cap 1{,}000 tokens & 1{,}000 / 1{,}000 / 1{,}000 & 0 / 0 / 0 \\
\midrule
\mGptFiveFive & B, C & 2026-04-24 & OpenAI & -- & effort medium & -- / 512 / 512 & -- / 0 / 7 \\
\mGlmFiveOne & B, C & 2026-04-07 & SiliconFlow & fp8 & cap 1{,}000 tokens & -- / 449 / 411 & -- / 0 / 0 \\
\mKimiKTwoSix & B, C & 2026-04-20 & SiliconFlow & fp8 & cap 1{,}000 tokens & -- / 1{,}000 / 1{,}000 & -- / 1 / 4 \\
\mDeepSeekVFourPro & B, C & 2026-04-24 & Alibaba & fp8 & cap 1{,}000 tokens & -- / 1{,}000 / 1{,}000 & -- / 9 / 17 \\
\bottomrule
\end{tabular}
\endgroup

\end{table*}

\section{The prompt}
\label{app:prompt}
Figure~\ref{fig:prompt} gives the template every scored call used, with the
placeholders in braces filled per row. There is no price, bid, ask, depth or book
anywhere in it. When a row has no news, the news block is the single line
\code{No news is available.} A reply containing no parseable probability is
re-asked once with the strict suffix shown beneath the template.

\begin{figure*}[!tbp]
\centering
\begin{minipage}{0.95\linewidth}
\footnotesize
\begin{verbatim}
You are forecasting the outcome of a prediction-market question.

Today's date and time (UTC): {today}
Everything below is what is known as of today. Do not use or assume
knowledge of anything that happened after today.

Event: {event_title}
Question: {question}{contract_line}

Resolution rules (word for word from the exchange):
{description}

Scheduled market close: {market_end_date}

News:
{news_block}

Give your own probability that this question resolves Yes. Reason
briefly, then end your reply with one line in exactly this form:
PROBABILITY: <number between 0 and 1>
\end{verbatim}
\rule{\linewidth}{0.4pt}
\begin{verbatim}
Your previous reply did not contain a usable probability. You must
answer with a number even if you are uncertain; refusing is not an
option. Reply with exactly one line and nothing else:
PROBABILITY: <decimal number between 0 and 1, for example 0.35>
\end{verbatim}
\end{minipage}
\caption{The forecasting prompt (above) and the strict re-ask appended after an unparseable
reply (below). Braces mark per-row fields. Line breaks inside the running sentences are added
here for column width; the template itself does not wrap them. Source:
\code{harness/prompt.py}.}
\label{fig:prompt}
\end{figure*}

\section{The memory probe}
\label{app:memory}
Every row carrying news was asked twice, once with the news and once with it
removed. The pair serves two purposes.

\textbf{Flagging memorisation.} A model that answers a news-bearing question just
as confidently and just as correctly without the news may be recalling the
outcome rather than reasoning from what it was given. Rows meeting that condition
are removed from the convergence test of Appendix~\ref{app:convergence}, the only
analysis the quarantine touches. Across the 21 probe runs (the 20 scored-model runs and the validation run) the flag
rate is 0.77 to 2.32 per cent.

\textbf{Measuring what news is worth.} The same pairing gives the value of news
directly, as the Brier score without it minus the Brier score with it. On the test split it ranges from $-0.001$ to $+0.017$ by model and window and is
positive in 18 of the 20 model-window cells, but small beside the gap to the
market.

\textbf{Validation, and why it is inconclusive.} To check that the probe detects
anything, a model published well after \winA settled (\mbox{GLM-5.3 Flash}, 26
August 2026) was run over 500 fixed \winA
rows. If the probe worked, that model should have been flagged more often than
the eligible models on the same rows. It was flagged less often: 1.4 per cent,
seven rows of 500, against 2.2 per cent for \mGptFiveTwo, 1.6 for
\mDeepSeekVThreeTwo, 2.2 for \mKimiKTwoFive and 3.4 for \mGptOss. Either the
probe lacks power, or that model had not memorised those particular markets. The
exercise cannot separate the two, and we do not claim the probe as a working
leakage detector; the release-date rule is what the leakage argument rests on.
